\documentclass[11pt]{article}
\usepackage[T1]{fontenc}
\usepackage{lmodern}
\usepackage{amsmath,amssymb,amsthm,mathtools}
\usepackage[margin=1in]{geometry}
\usepackage{booktabs,microtype,hyperref}
\hypersetup{colorlinks=true,linkcolor=blue,citecolor=blue,urlcolor=blue,
pdftitle={Explicit constructions and global bounds for the Mordell discriminant problem},
pdfauthor={Fengcan Jiang}}
\setlength{\emergencystretch}{2em}
\newtheorem{theorem}{Theorem}[section]
\newtheorem{proposition}[theorem]{Proposition}
\newtheorem{lemma}[theorem]{Lemma}
\newtheorem{corollary}[theorem]{Corollary}
\theoremstyle{remark}
\newtheorem{remark}[theorem]{Remark}
\newcommand{\C}{\mathbb C}
\newcommand{\R}{\mathbb R}
\newcommand{\Q}{\mathbb Q}
\newcommand{\E}{\mathbb E}
\numberwithin{equation}{section}
\title{Explicit constructions and global bounds\\for the Mordell discriminant problem}
\author{Fengcan Jiang\thanks{Email:
\href{mailto:2255214397@stu.xjtu.edu.cn}{\texttt{2255214397@stu.xjtu.edu.cn}}.}}
\date{October 3, 2026}
\begin{document}
\maketitle
\begin{abstract}
This research note studies the largest squared Vandermonde product
under a prescribed quadratic moment. Elementary Gaussian averages
and Hadamard--Jensen estimates give explicit bounds for every point
count, determining the maximum to within a factor $\exp(O(n))$.
They recover its leading quadratic and $n\log n$ asymptotic terms.
A deterministic lattice construction attains the leading quadratic
term and has an exactly computable rational product. Resultants give
explicit multi-ring constructions; a ten-point example exceeds the
center-plus-nine-gon competitor. For six points we prove the sharp
maximum among configurations containing the origin, and exclude two
concentric equilateral triangles. An explicit neighborhood certificate
controls the central pentagon. The central nine-gon is a strict local
maximum despite being globally suboptimal, while an explicit Fourier
perturbation improves the central polygon for every $n\ge11$.
A certified grid error bound gives arbitrary global accuracy.
A polynomial trace test enforces genuine conjugate pairing, and
low-degree critical equations yield a finite terminating exact-value
algorithm, demonstrated for two and three points. These results do
not determine the unrestricted six-point maximum or a closed formula
for every finite point count.
\end{abstract}

\section{Scope and normalization}
For $n\ge2$, put
\[
 \Delta_n(z)=\prod_{i<j}|z_i-z_j|^2,\qquad
 M_n=\max_{\sum_i|z_i|^2=n}\Delta_n(z),\qquad
 N=\frac{n(n-1)}2,\qquad P_n=\prod_{k=0}^{n-1}k!.
\]
The maximum exists by compactness; it is positive and has no
collisions, since a unit regular polygon has value $n^n$.
Every maximizer has zero centroid: centering preserves the distances
and decreases the quadratic moment by $n$ times the squared centroid;
rescaling a nonzero centroid configuration to the constraint would
strictly increase its product.

The five-point equality $M_5=5^5$ is proved by Wang \cite{Wang}.
Chalk's center-plus-polygon construction \cite{Chalk} disproves
$M_n=n^n$ for $n\ge6$. The present note provides general constructions
and global bounds, rather than a claimed formula for every finite $M_n$.
The asymptotic context is classical Gaussian weighted Fekete theory
and the planar logarithmic gas \cite{AOC,SS}. No claim of priority is
made for these asymptotic terms or elementary bounds.

\section{Explicit universal bounds}
\label{general_sec_bounds}
\begin{theorem}\label{general_thm_bounds}
Put $S_m(x)=\sum_{k=0}^m x^k/k!$. For every $\alpha>0$,
\begin{equation}\label{general_eq_bounds}
 L_n:=\frac{n!(n-2)!\,n^N P_n}{(N+n-2)!}
 \le M_n\le U_n(\alpha):=\frac{P_n S_{n-1}(\alpha)^n}{\alpha^N}.
\end{equation}
The upper bound has a unique minimizing parameter $\alpha_n>0$,
characterized by
\begin{equation}\label{general_eq_optimal_alpha}
 \sum_{k=0}^{n-1}(2k-(n-1))\frac{\alpha_n^k}{k!}=0.
\end{equation}
In particular,
\begin{equation}\label{general_eq_exponential_upper}
 M_n\le\widehat U_n:=e^N\left(\frac nN\right)^NP_n.
\end{equation}
\end{theorem}
\begin{proof}
Take independent standard complex Gaussian variables $Z_j$,
each with density $\pi^{-1}e^{-|z|^2}$. Their moments satisfy
$\E(Z^k\overline Z^{\,\ell})=0$ for $k\ne\ell$ and $k!$ for $k=\ell$.
Expansion of the squared Vandermonde determinant gives
$\E\Delta_n(Z)=n!P_n$.

Let $Y_j=Z_j-n^{-1}\sum_iZ_i$ and $R^2=\sum_j|Y_j|^2$.
A unitary change of Gaussian coordinates shows that $R^2$ has the
$\Gamma(n-1,1)$ distribution and is independent of
$W=\sqrt n\,Y/R$, the uniform direction on the zero-centroid sphere.
Translation invariance and homogeneity give
\[
 n!P_n=\E\Delta_n(Y)
 =\frac{\Gamma(N+n-1)}{\Gamma(n-1)n^N}\,\E\Delta_n(W).
\]
Thus $L_n=\E\Delta_n(W)\le M_n$.

For the upper bound, apply row Hadamard to
$A_{jk}=\alpha^{k/2}z_j^k/\sqrt{k!}$, $0\le k<n$.
Writing $t_j=|z_j|^2$, this gives
\[
 \frac{\alpha^N}{P_n}\Delta_n(z)
 \le\prod_{j=1}^nS_{n-1}(\alpha t_j)
 \le S_{n-1}(\alpha)^n.
\]
The second inequality is Jensen, since $\sum_jt_j=n$
and $\log S_m$ is concave on $[0,\infty)$.
For $m=1$ this is immediate. For $m\ge2$, use $S'_m=S_{m-1}$,
$S''_m=S_{m-2}$ and
\[
 S_{m-1}(x)^2-S_m(x)S_{m-2}(x)
 =\frac{x^{m-1}}{(m-1)!}
 \left(1+\sum_{k=1}^{m-1}
         \left(1-\frac km\right)\frac{x^k}{k!}\right)\ge0.
\]
This proves \eqref{general_eq_bounds}.

Set $m=n-1$ and $q(\alpha)=\alpha S'_m(\alpha)/S_m(\alpha)$.
For the distribution on $\{0,\ldots,m\}$ with probabilities
proportional to $\alpha^k/k!$, $q$ is the mean of $k$, and
$dq/d\log\alpha$ is its strictly positive variance.
Thus $q$ increases from $0$ to $m$. The logarithmic derivative of
$U_n(\alpha)$ with respect to $\log\alpha$ is $nq(\alpha)-N$,
so its unique zero is the global minimum, giving
\eqref{general_eq_optimal_alpha}. Finally, $S_m(\alpha)\le e^\alpha$
and $\alpha=N/n$ prove \eqref{general_eq_exponential_upper}.
\end{proof}

The lower bound is an average, not a proposed extremal configuration.
For rational $\alpha$, $U_n(\alpha)$ is rational, giving a fully exact
finite interval. The optimal parameter is algebraic by
\eqref{general_eq_optimal_alpha} and can be enclosed by polynomial
root isolation. For $n=2$, $L_2=U_2(1)=4$.
For example, $\alpha=11/4$ gives the exact bound $M_6<89309$, whereas
the central pentagon gives $M_6\ge48146.9424205824$.
The decimal lower endpoint here is exact; its value is $6^{15}/5^{10}$.

\begin{corollary}\label{general_cor_asymptotic}
With $c=\tfrac12\log2-\tfrac14$,
\begin{equation}\label{general_eq_asymptotic}
 \log M_n=cn^2+\frac n2\log n+O(n).
\end{equation}
\end{corollary}
\begin{proof}
Summing Stirling's formula, or applying Euler--Maclaurin to
$\log P_n=n\log\Gamma(n)-\sum_{k=1}^{n-1}k\log k$, gives
\[
 \log P_n=\frac{n^2}{2}\log n-\frac{3n^2}{4}
            +\frac n2\log(2\pi)+O(\log n).
\]
Substitution in \eqref{general_eq_exponential_upper} yields
\[
 \log\widehat U_n=cn^2+\frac n2\log n
                      +\frac n2\log\pi+O(\log n).
\]
The exact ratio of the upper and lower bounds is
\[
 \frac{\widehat U_n}{L_n}
 =e^NN^{-N}\frac{(N+n-2)!}{n!(n-2)!}.
\]
Stirling gives the logarithm of this ratio as
$(2-\log2)n+O(\log n)$, proving the assertion.
\end{proof}

\section{A deterministic construction for every number of points}
\label{general_sec_lattice}
The uniform disk limit is a classical feature of quadratic weighted
Fekete problems; see \cite{AOC,SS}. The following construction supplies
explicit points and exact finite products. Its optimality statement
concerns the leading $n^2$ scale, not the maximum at each finite $n$.

\begin{theorem}\label{general_thm_lattice}
Order the Gaussian integers $v\in\mathbb Z+i\mathbb Z$ by increasing
$|v|^2$, breaking ties lexicographically by $(\Re v,\Im v)$, and take
the first $n$ distinct points $v_1,\ldots,v_n$. Define
\begin{equation}\label{general_eq_lattice}
 \mu_v=\frac1n\sum_jv_j,\quad
 D_n=n\sum_j|v_j|^2-\left|\sum_jv_j\right|^2,\quad
 w_j=\frac{n(v_j-\mu_v)}{\sqrt{D_n}}.
\end{equation}
Then $D_n$ is a positive integer, $\sum_jw_j=0$,
$\sum_j|w_j|^2=n$, and
\begin{equation}\label{general_eq_lattice_product}
 \Delta_n(w)=\left(\frac{n^2}{D_n}\right)^N
                 \prod_{i<j}|v_i-v_j|^2\in\Q_{>0}.
\end{equation}
Furthermore,
\begin{equation}\label{general_eq_lattice_asymptotic}
 \lim_{n\to\infty}\frac{\log\Delta_n(w)}{n^2}=c,\qquad
 \log M_n-\log\Delta_n(w)=o(n^2).
\end{equation}
\end{theorem}
\begin{proof}
Centering and scaling give the finite identities. Positivity of $D_n$
follows because the points are distinct and $n\ge2$.
Put $R_n=\max_{j\le n}|v_j|$ and $a_n=n/\sqrt{D_n}$.
Covering lattice points by unit squares proves
\[
 \#\{v:|v|\le R\}=\pi R^2+O(R+1),\qquad
 \sum_{|v|\le R}|v|^2=\frac\pi2R^4+O(R^3+1).
\]
Only $O(R_n+1)$ boundary points can be omitted when a radius shell
is cut by the tie rule. The full disk has zero coordinate sum, so
$\sum_{j\le n}v_j=O(R_n^2+1)$ and $\mu_v=O(1)$.
Consequently
\[
 R_n^2\sim\frac n\pi,\quad
 \sum_j|v_j-\mu_v|^2\sim\frac{n^2}{2\pi},\quad
 a_n^2\sim\frac{2\pi}{n},\quad a_nR_n\longrightarrow\sqrt2.
\]
Riemann sums over unit squares show that
$\mu_n=n^{-1}\sum_j\delta_{w_j}$ converges weakly to the uniform
probability measure $\mu$ on $|z|\le\sqrt2$.
Their supports lie in a common bounded disk for large $n$.

The minimum pair distance is at least $a_n$.
Packing lattice squares gives, for $r\ge a_n$, at most
$C(r/a_n)^2\le C'nr^2$ other points within distance $r$ of each point.
Thus the ordered-pair count $H_n(r)$ is zero for $r<a_n$ and
satisfies $H_n(r)\le C'n^2r^2$ for $r\ge a_n$.
For $0<\varepsilon<1$, layer integration gives
\begin{align*}
 \frac1{n^2}\sum_{\substack{i\ne j\\|w_i-w_j|<\varepsilon}}
       -\log|w_i-w_j|
 &=\frac{H_n(\varepsilon)}{n^2}|\log\varepsilon|
       +\frac1{n^2}\int_0^\varepsilon\frac{H_n(r)}r\,dr\\
 &\le C'\varepsilon^2
             \left(|\log\varepsilon|+\frac12\right).
\end{align*}
For fixed $\varepsilon$, the truncated kernel
$\max\{\log|z-u|,\log\varepsilon\}$ is continuous on the common
compact support. Weak convergence gives convergence of its double
integrals, and its diagonal empirical contribution is
$\log\varepsilon/n$.
First let $n\to\infty$ and then $\varepsilon\downarrow0$ to obtain
\[
 \frac{\log\Delta_n(w)}{n^2}
 =\frac1{n^2}\sum_{i\ne j}\log|w_i-w_j|
 \longrightarrow\iint\log|z-u|\,d\mu(z)d\mu(u).
\]
For $R=\sqrt2$, the angular mean identity
\[
 \frac1{2\pi}\int_0^{2\pi}\log|r-se^{i\theta}|\,d\theta
 =\log\max\{r,s\}
\]
gives the disk potential
$\log R-\tfrac12+|z|^2/(2R^2)$ for $|z|\le R$.
Since $\int|z|^2d\mu=R^2/2$, its average is $\log R-1/4=c$.
This proves the limit. The remaining assertion follows from
Corollary~\ref{general_cor_asymptotic} and $\Delta_n(w)\le M_n$.
\end{proof}

\section{Exact multi-ring lower-bound constructions}
Let $\delta\in\{0,1\}$ indicate the presence of the origin.
Ring $a$ consists of $m_a\ge2$ points
$r_a\exp(i(\theta_a+2\pi j/m_a))$, $0\le j<m_a$, with $r_a>0$.
Require
\[
 n=\delta+\sum_a m_a,\qquad \sum_a m_ar_a^2=n.
\]
For two rings define $g_{ab}=\gcd(m_a,m_b)$ and
$L_{ab}=\operatorname{lcm}(m_a,m_b)$.

\begin{proposition}\label{prop:rings}
The exact product of this configuration is
\begin{equation}\label{eq:rings}
\begin{split}
 \Delta_n={}&\prod_a m_a^{m_a}
       r_a^{m_a(m_a-1)+2\delta m_a}\\
 &\quad\cdot\prod_{a<b}
 \left|r_a^{L_{ab}}e^{iL_{ab}\theta_a}
       -r_b^{L_{ab}}e^{iL_{ab}\theta_b}\right|^{2g_{ab}}.
\end{split}
\end{equation}
Consequently this expression is a lower bound for $M_n$.
\end{proposition}
\begin{proof}
Each ring is the root set of
$z^{m_a}-r_a^{m_a}e^{im_a\theta_a}$, whose absolute discriminant
is $m_a^{m_a}r_a^{m_a(m_a-1)}$.
The absolute resultant of $z^m-A$ and $z^k-B$ is
$|A^{k/g}-B^{m/g}|^g$, where $g=\gcd(m,k)$:
the values of $z^k$ on the first root set traverse the roots of
$u^{m/g}-A^{k/g}$, each with multiplicity $g$.
Squaring this resultant gives every cross-ring squared distance.
Each origin-to-ring product contributes $r_a^{2m_a}$.
Multiplying proves the formula, also when cross-ring collisions make
both sides zero.
\end{proof}

With two rings of populations $a,b$ and $\delta$ origin points,
write $g=\gcd(a,b)$, $L=ab/g$, and $x=r^2/s^2$.
The optimal relative phase within this family is
$\theta_2-\theta_1=\pi/L$ modulo $2\pi/L$.
Set
$A=a(a-1)/2+\delta a$ and $B=b(b-1)/2+\delta b$.
Then $N=A+B+ab$, $s^2=n/(ax+b)$, and the family maximum reduces to
\begin{equation}\label{eq:two-rings}
 \max_{x>0} a^ab^bn^N
       \frac{x^A(1+x^{L/2})^{2g}}{(ax+b)^N}.
\end{equation}
The interior critical equation is
\begin{equation}\label{eq:ring-critical}
 Ab-a(B+ab)x+b(A+ab)x^{L/2}-aBx^{L/2+1}=0.
\end{equation}
It follows by logarithmic differentiation and
$gL=ab$. For odd $L$, substituting $x=\rho^2$ makes this a polynomial
with integer exponents, so exact root isolation and value comparison
can certify the optimum within the chosen family.
For three or more rings, the relative phases must be optimized together;
pairwise opposite phases need not be simultaneously possible.
Family optimization alone does not prove unrestricted optimality.

\subsection{A strict ten-point improvement on the central polygon}
The center-plus-$(n-1)$-gon value is
\[
 C_n=(n-1)^{n-1}\left(\frac n{n-1}\right)^N.
\]
For $n=10$, consider instead the roots of
\begin{equation}\label{eq:ten-polynomial}
 (z^2-u)(z^8+v^4),\qquad u=\frac18,\quad v=\frac{39}{32}.
\end{equation}
They are $\pm\sqrt u$ and the vertices of a regular octagon of
radius $\sqrt v$, rotated by $\pi/8$.
They satisfy $2u+8v=10$ and have product
\begin{equation}\label{eq:ten-product}
 D_{10}=4u\,8^8v^{28}(u^4+v^4)^4,\qquad
 \frac{D_{10}}{C_{10}}>\frac{113}{100}>1.
\end{equation}
The last inequality is a rational comparison verified by the
accompanying program. Thus the central polygon cannot be the general
finite-$n$ answer. This two-ring example is a certified lower bound,
and is not asserted to maximize $\Delta_{10}$.

\section{Exact six-point subclasses and a majorant obstruction}
Write $C_6=6^{15}/5^{10}$, the value of the origin together with a
regular pentagon of radius $\sqrt{6/5}$. The following results do
not establish $M_6=C_6$.

\begin{proposition}[Configurations containing the origin]
If $\sum_{i=1}^6|z_i|^2=6$ and some $z_i=0$, then $\Delta_6\le C_6$.
Equality holds precisely for the origin and a regular pentagon of
radius $\sqrt{6/5}$, up to rotation and permutation.
\end{proposition}
\begin{proof}
Assume $z_6=0$. Apply the five-point theorem of Wang \cite{Wang} to
$\sqrt{5/6}(z_1,\ldots,z_5)$, and arithmetic--geometric mean to
the five squared radii. Homogeneity gives
\[
 \Delta_6=\Delta_5(z_1,\ldots,z_5)\prod_{i=1}^5|z_i|^2
 \le5^5(6/5)^{10}(6/5)^5=C_6.
\]
The equality conditions give exactly the asserted configuration.
\end{proof}

\begin{proposition}[Two concentric equilateral triangles]
Among configurations consisting of two equilateral triangles centered
at the origin with $\sum_i|z_i|^2=6$, the largest product is $6^6$,
attained precisely by a unit regular hexagon.
\end{proposition}
\begin{proof}
Let $a,b\ge0$ be the squared radii, so $a+b=2$.
If $ab=0$ there are collisions. Otherwise the ring discriminants are
$3^3a^3$ and $3^3b^3$. Their resultant gives the cross-distance product
$|a^{3/2}-b^{3/2}e^{3i\phi}|^6$, where $\phi$ is the relative rotation.
Thus, with $p=\sqrt{ab}\in(0,1]$,
\[
 \Delta_6\le3^6(ab)^3(a^{3/2}+b^{3/2})^6
 =3^6\{8p^2-6p^4+2p^5\}^3.
\]
The derivative of the expression in braces is
$2p(8-12p^2+5p^3)>0$: the last factor decreases on $[0,1]$
and equals $1$ at $p=1$. Its maximum is $4$, hence
$\Delta_6\le3^6\cdot4^3=6^6$.
Equality requires $a=b=1$ and $3\phi=\pi$ modulo $2\pi$,
giving a regular hexagon. Finally $6^6<C_6$ because
$C_6/6^6=6^9/5^{10}>1$.
\end{proof}

\subsection{Hermite majorants adapted to a candidate}
Suppose a candidate has $r$ distinct positive squared distances
$a_1,\ldots,a_r$. Let $H$ be the Hermite polynomial of degree at
most $2r-1$ satisfying $H(a_j)=\log a_j$ and $H'(a_j)=1/a_j$.
Then
\[
 \log x\le H(x)\quad(x>0),\qquad
 \log\Delta_n\le\sum_{i<j}H(|z_i-z_j|^2).
\]
The interpolation remainder is a nonpositive multiple of
$P(x)=\prod_j(x-a_j)^2$, since
$(\log x)^{(2r)}=-(2r-1)!/x^{2r}<0$.
The sum of $H$ is a finite combination of distance moments;
stationary identities and second-order Gram matrices supply
necessary constraints on those moments.

If an upper endpoint $R>\max_j a_j$ is available, set
$\kappa_R=(H(R)-\log R)/P(R)>0$.
Adding the interpolation node $R$ gives
$\log x\le H(x)-\kappa_RP(x)$ for $0<x\le R$:
the next derivative is positive and the remainder has the sign of
$P(x)(x-R)$. Any verified $0<k\le\kappa_R$ also gives a valid majorant.

For the central pentagon the nodes are
\[
 a_1=6/5,\qquad
 a_2=3(5-\sqrt5)/5,\qquad a_3=3(5+\sqrt5)/5,
\]
each occurring five times.
A regular hexagon has squared-distance multiplicities $6,6,3$ at
$1,3,4$, respectively. The exact interval check
\texttt{verify\_hermite.py} shows that its degree-five majorant satisfies
\[
 \frac2{25}<6H(1)+6H(3)+3H(4)-\log C_6<\frac{83}{1000}.
\]
Even the maximal interpolation correction at $R=15/2$ leaves
this gap greater than $3/100$.
Thus this majorant alone is insufficient. A sharp six-point argument
still needs further moment information or a separate comparison.
The subclass results above do not settle the unrestricted optimum.

\section{Local certificates and explicit improving perturbations}
\label{sec:local-certificates}
The spectral stability of polygonal vortex equilibria has a substantial
classical literature; see Cabral--Schmidt \cite{CS}.
We give direct constrained-discriminant calculations, with explicit
Taylor bounds and an exact treatment of a degenerate ten-point case.
We make no claim of priority for the polygonal stability thresholds.

\subsection{The complete central-polygon spectrum}
Write $n=m+1$, $m\ge4$, $r^2=n/m$, $\omega=e^{2\pi i/m}$, and
\[
 z^*=(0,r,r\omega,\ldots,r\omega^{m-1}),\qquad \lambda=m/2.
\]
The root polynomial $z(z^m-r^m)$ or a geometric sum verifies
\[
 \sum_{j\ne i}\frac1{z_i^*-z_j^*}=\lambda\overline{z_i^*}.
\]
Define the complex symmetric matrix
\[
 B_{ij}=-\frac1{(z_i^*-z_j^*)^2}\quad(i\ne j),\qquad
 B_{ii}=\sum_{j\ne i}\frac1{(z_i^*-z_j^*)^2}.
\]
Let $\langle x,y\rangle=\sum_i\overline{x_i}y_i$, and put
\[
 E=\{h\in\C^n:\ \sum_i h_i=0,\ \langle z^*,h\rangle=0\}.
\]
This is a local slice for the zero-centroid sphere modulo rotation.
Its normalized chart is
\begin{equation}\label{eq:local-chart}
 z(h)=\frac{\sqrt n\,(z^*+h)}{\sqrt{n+\|h\|^2}}.
\end{equation}
The Taylor expansion of
$F(h)=\log\Delta_n(z(h))-\log C_n$ has quadratic term
\begin{equation}\label{eq:quadratic-slice}
 Q_2(h)=-\Re(h^{\mathsf T}Bh)-\lambda\|h\|^2.
\end{equation}
Indeed, expand each $2\Re\log(1+(h_i-h_j)/(z_i^*-z_j^*))$
and subtract $N\log(1+\|h\|^2/n)$; the linear term vanishes
by stationarity and $h\in E$.

Set
\[
 f_k=\frac1{\sqrt m}(0,1,\omega^k,\ldots,\omega^{(m-1)k}),
 \quad 1\le k<m,\qquad
 u=\frac{(-m,1,\ldots,1)}{\sqrt{m(m+1)}}.
\]
Together with $\mathbf1/\sqrt n$, these form a Hermitian orthonormal
basis. Direct divided-difference expansion and the identity
$\sum_{j=0}^{m-1}\omega^{jk}=0$ for $m\nmid k$ give
\begin{align}
 B\mathbf1&=0,& Bf_1&=\lambda f_{m-1},\notag\\
 Bf_2&=\gamma u,& Bu&=\gamma f_{m-2},\label{eq:central-spectrum-map}\\
 Bf_k&=c_kf_{k-2}\quad(3\le k<m),&
 \gamma&=\frac m{\sqrt{m+1}},\quad
 c_k=\frac{1-(k-2)(m-k)/2}{r^2}.\notag
\end{align}
For example, on ring coordinates and $3\le k<m$, the expansion is
\[
 (B(z^*)^{\circ k})_i
 =\bigl(1-(k-2)(m-k)/2\bigr)(z_i^*)^{k-2};
\]
the center coordinate is zero, since its corresponding power sum
vanishes. For $k=2$, the center coordinate supplies the vector $u$.
The output vectors in \eqref{eq:central-spectrum-map} are distinct
orthonormal basis vectors, so the singular values of $B$ are
\begin{equation}\label{eq:central-singular-values}
 0,\quad\lambda,\quad\gamma,\quad\gamma,
 \quad |c_k|\ (3\le k<m).
\end{equation}
On $E$ the singular value $\lambda$ associated with $f_1$
has been removed. Since
\[
 \max_{3\le k<m}(k-2)(m-k)
 =\left\lfloor\frac{(m-2)^2}{4}\right\rfloor,
\]
the nontrivial singular values are strictly below $\lambda$
for $4\le m\le8$. Thus the central polygon is a strict local
maximum modulo rotations for $5\le n\le9$.
For $m=9$, $c_5=c_6=-\lambda$; the remaining singular values on $E$
are strictly smaller. For $m\ge10$, some $|c_k|>\lambda$.
An explicit improving direction for that case is given below.

\subsection{An explicit six-point neighborhood certificate}
\begin{theorem}\label{thm:six-basin}
Take $n=6$ and $z^*$ as above. For $h\in E$ with
$0<\|h\|\le1/25$, the chart \eqref{eq:local-chart} satisfies
\begin{equation}\label{eq:six-basin}
 \log\Delta_6(z(h))\le\log C_6-\frac14\|h\|^2.
\end{equation}
\end{theorem}
\begin{proof}
For $m=5$, $c_3=c_4=0$, and the norm of $B$ restricted to $E$
is $\gamma=5/\sqrt6$. Hence \eqref{eq:quadratic-slice} gives
$Q_2(h)\le-\delta\|h\|^2$, where
$\delta=5/2-5/\sqrt6>9/20$.
The last strict inequality follows from $25/6<(41/20)^2$.

Write $\rho=\|h\|$ and
$v_{ij}=(h_i-h_j)/(z_i^*-z_j^*)$.
The smallest squared distance at $z^*$ is $r^2=6/5$, so
\[
 \max_{i<j}|v_{ij}|\le\sqrt{5/3}\,\rho<\frac43\rho=:b_\rho,
 \qquad
 \sum_{i<j}|v_{ij}|^2\le\frac6{r^2}\rho^2=5\rho^2.
\]
Here the second inequality uses
$\sum_{i<j}|h_i-h_j|^2=6\|h\|^2$ for $\sum_i h_i=0$.
For $|v|<1$,
\[
 |\log(1+v)-v+v^2/2|
 \le\frac{|v|^3}{3(1-|v|)}.
\]
Also $\log(1+x)\ge x-x^2/2$ for $x\ge0$.
Applying these bounds to the exact chart product gives
\[
 F(h)\le-\delta\rho^2+
 \left(\frac{10b_\rho}{3(1-b_\rho)}+\frac5{24}\rho^2\right)\rho^2.
\]
For $\rho\le1/25$, the coefficient in parentheses is at most
$40/213+1/3000<1/5$. Consequently
$F(h)<-(9/20-1/5)\rho^2=-\rho^2/4$.
This proves the stated non-strict bound.
\end{proof}

This theorem excludes a concrete neighborhood from a search for a
strictly better six-point competitor. It makes no comparison with
configurations outside that neighborhood.

\subsection{Ten points: a strict local maximum that is not global}
\begin{theorem}\label{thm:ten-local-trap}
For $n=10$, the origin plus a regular nine-gon of radius $\sqrt{10/9}$
is a strict local maximum modulo rotations. It is not a global maximum.
\end{theorem}
\begin{proof}
The failure of global optimality is already certified by
\eqref{eq:ten-product}. We prove strict local maximality.
Now $m=9$, $r=\sqrt{10/9}$ and $\omega=e^{2\pi i/9}$.
The only null directions of $Q_2$ on $E$ are
\[
 h(a)_0=0,\qquad
 h(a)_{j+1}=r(a\omega^{5j}+\overline a\,\omega^{6j}),
 \qquad a\in\C.
\]
In fact, on the two relevant Fourier modes, writing their
coefficients as $a,b$ gives $Q_2=-45|a-\overline b|^2$.
Let $W\subset E$ be the real orthogonal complement of this
two-dimensional nullspace. On $W$ write
$Q_2(v)=-(v,Av)_{\R}$, with $A$ positive definite.

For $h=h(1)$ put
$u_{ij}=(h_i-h_j)/(z_i^*-z_j^*)$ and
$a_i=\sum_{j\ne i}u_{ij}^2/(z_i^*-z_j^*)$.
Finite root-of-unity sums give
\begin{align*}
 \sum_{i<j}u_{ij}&=0,&\sum_{i<j}u_{ij}^2&=-90,\\
 \sum_{i<j}u_{ij}^3&=0,&\sum_{i<j}u_{ij}^4&=2970,\\
 a_0&=-9/r,&
 a_{j+1}&=(1-15\omega^{7j}-10\omega^{8j})/r.
\end{align*}
These finite identities are also checked exactly in
$\Q(\omega)$ with $\omega^6+\omega^3+1=0$ by
\texttt{verify\_polygon\_\allowbreak perturbations.py}.
The first four sums can independently be obtained by expanding the
divided differences and using
\[
 \sum_{i\ne j}\omega^{pi+qj}
 =m^2\mathbf1_{m\mid p}\mathbf1_{m\mid q}
       -m\mathbf1_{m\mid p+q}.
\]

The straight null-direction quartic term is
$-2970/2+45\cdot2^2/2=-1395$.
For a transverse correction $th+t^2k$, $k\in W$, the coefficient
of $t^4$ is
\begin{equation}\label{eq:ten-effective-quartic}
 -1395+Q_2(k)+L(k),\qquad
 L(k)=2\Re\sum_i a_i k_i.
\end{equation}
The optimizing correction is
\[
 (k_0)_0=-12r,\qquad
 (k_0)_{j+1}=r(4/3-(17/3)\omega^{2j}).
\]
It lies in $W$. The projected real gradient of $L$ on $E$ is
$2\overline a+18z^*$, and direct substitution in
\eqref{eq:central-spectrum-map} gives
$2Ak_0=2\overline a+18z^*$.
Positive definiteness on $W$ therefore makes $k_0$ the unique maximizer
of $Q_2(k)+L(k)$ there. Writing $k=r(-9c,c+d\omega^{2j})$ reduces the equation for the
correction to $5c+\overline d=1$, $5d+10\overline c=-15$, hence
$c=4/3$, $d=-17/3$. Fourier orthogonality gives
$\|k\|^2=r^2(90c^2+9d^2)$,
$k^TBk=180cd$, and $L(k)=180c-270d$ at this real solution.
The exact constants are
\[
 \|k_0\|^2=\frac{4490}{9},\quad
 \Re(k_0^{\mathsf T}Bk_0)=-1360,\quad
 Q_2(k_0)=-885,\quad L(k_0)=1770.
\]
Thus the largest effective quartic coefficient is
$-1395+885=-510<0$.

For completeness, the negative quartic coefficient controls every
null direction, rather than only $h(1)$.
The cyclic relabeling combined with physical rotation that fixes
$z^*$ acts by $a\mapsto\omega^{-4}a$ on the nullspace.
As this phase has order nine, its invariant homogeneous cubic
polynomial is zero and its invariant homogeneous quartic polynomial
is a multiple of $|a|^4$.
The implicit function theorem on $W$ gives a unique nearby transverse
maximizer $w(a)=O(|a|^2)$ for each small $a$.
Its reduced value is therefore
\[
 F(h(a)+w(a))=-510|a|^4+O(|a|^5).
\]
It is strictly negative for small $a\ne0$.
Continuity preserves strict concavity in the transverse variables,
so any other nearby transverse choice has smaller value.
This proves strict local maximality on the normalized slice.
Nonzero centroids strictly decrease the normalized value relative
to their centered configuration, so the conclusion also holds on
the full constraint, modulo rotations.
\end{proof}

This exact local trap explains why even careful local searches and
second-order tests cannot alone certify the finite global maximum.

\subsection{A certified improvement for every \texorpdfstring{$n\ge11$}{n >= 11}}
\begin{theorem}\label{thm:all-n-improvement}
Let $m=n-1\ge10$, $k=\lfloor(m+2)/2\rfloor$,
$q=\lfloor(m-2)^2/4\rfloor$, and
\[
 b=\begin{cases}1&m\text{ even},\\2&m\text{ odd},\end{cases}
 \qquad A_m=\frac{bm}{2}(q-m-3)>0,\qquad H=m+2.
\]
Choose the explicit rational step
\[
 t_m=\min\left\{\frac1{2H},\frac{A_m}{4NH^3}\right\}.
\]
Set $h_0=0$ and, on the ring,
\[
 h_{j+1}=\begin{cases}
 r\omega^{kj},&m\text{ even},\\
 r(\omega^{kj}+\omega^{(m+2-k)j}),&m\text{ odd}.
 \end{cases}
\]
Then the explicit configuration
\[
 y_j=\frac{z_j^*+t_mh_j}{\sqrt{1+bt_m^2}}
\]
has zero centroid, quadratic moment $n$, distinct points, and
\begin{equation}\label{eq:all-n-strict-gain}
 \log\Delta_n(y)-\log C_n\ge\frac23 A_m t_m^2>0.
\end{equation}
Consequently $M_n>C_n$ for every $n\ge10$, using the separate
ten-point construction when $n=10$.
\end{theorem}
\begin{proof}
Fourier orthogonality gives $\sum h_j=0$,
$\langle z^*,h\rangle=0$, and $\|h\|^2=bn$.
For even $m$, $2k-2=m$, and \eqref{eq:central-spectrum-map} gives
$h^{\mathsf T}Bh=m(1-q/2)$.
For odd $m$, the two modes pair and this value is doubled.
Thus the quadratic coefficient along the normalized direction is
\[
 -\Re(h^{\mathsf T}Bh)-Nb=A_m.
\]
To see its positivity, for even $m=2v\ge10$ one has
$q-m-3=v^2-4v-2>0$; for odd $m=2v+1\ge11$ it is
$(v-4)(v+1)>0$.

Put $u_{ij}=(h_i-h_j)/(z_i^*-z_j^*)$.
Each ring divided difference is a sum of $k$ unit-modulus terms,
or two such sums with total $m+2$ terms. Center-to-ring ratios
are bounded by $b$. Hence $|u_{ij}|\le H$.
The chosen step gives $Ht_m\le1/2$, so every factor
$1+t_mu_{ij}$ is nonzero.
The logarithmic Taylor bound and $\log(1+bt^2)\le bt^2$ give
\[
 \log\Delta_n(y)-\log C_n
 \ge A_mt_m^2-\frac{2NH^3t_m^3}{3(1-Ht_m)}
 \ge A_mt_m^2-\frac43NH^3t_m^3
 \ge\frac23 A_mt_m^2>0.
\]
This also verifies distinctness and proves the assertion.
\end{proof}

The gain bound is deliberately conservative. Its role is to certify
one explicit strictly improving configuration for every point count
in this range, rather than to approximate the unrestricted maximum.

\section{Finite-point numerical candidates}
The optional program \texttt{explore\_candidates.py} uses NumPy,
analytic gradients, BFGS updates and an Armijo line search to minimize
\[
 \frac{n-1}{2}\sum_i|z_i|^2-\log\Delta_n(z).
\]
Its stationary scale is $\sum_i|z_i|^2=n$.
For each $n=6,\ldots,12$, it uses 50 Gaussian random starts and
one center-plus-polygon start, with seed 20261004.
The table shows the largest observed logarithmic product.
These floating-point searches do not certify global optimality.

\begin{center}
\begin{tabular}{rrr}
\toprule
$n$ & Best observed $\log\Delta_n$ & Log gain over $C_n$\\
\midrule
6 & 10.7820129141 & $<10^{-12}$\\
7 & 13.9877210917 & $<10^{-12}$\\
8 & 17.3602500369 & $<10^{-12}$\\
9 & 20.8757216171 & $<10^{-12}$\\
10 & 24.7089672032 & 0.192722803\\
11 & 28.7906758481 & 0.522765029\\
12 & 33.1914966077 & 1.071897730\\
\bottomrule
\end{tabular}
\end{center}

Each best coordinate set is also rounded to Gaussian integers on
a grid of spacing $10^{-6}$, then centered and normalized using
\eqref{general_eq_lattice}, without any lattice-selection assumption.
The exact product formula \eqref{general_eq_lattice_product} applies
to any distinct Gaussian integers. The corresponding rationalized
ten-point candidate strictly exceeds $(121/100)C_{10}$.
Its integer coordinates, normalization denominator and exact rational
product are saved in \texttt{numerical\_candidates.json}.
This last comparison is an exact attainable lower bound; it does
not convert the numerical search into a global-optimality proof.
Gradient residuals and convergence flags are also recorded, including
searches that did not reach the requested gradient tolerance.


\section{Certified approximation of the global maximum}
\label{sec:grid-global}
A grid maximum becomes useful for the unrestricted problem only when
its discretization error is bounded uniformly. At an actual maximum,
stationarity cancels the first-order rounding error.

\begin{lemma}\label{lem:global-separation}
Every maximizer satisfies
\begin{equation}\label{eq:global-separation}
 \sum_{j\ne i}\frac1{z_i-z_j}=\lambda\overline z_i,
 \qquad \lambda=\frac{n-1}{2},\qquad
 |z_i-z_j|^2\ge\frac{2n}{(n-1)^2}.
\end{equation}
\end{lemma}
\begin{proof}
The multiplier follows by multiplying the stationary equations by
$z_i$ and summing, since each pair contributes one.
Homogeneity shows that a constrained global maximizer also maximizes
$\log\Delta_n(z)-\lambda\sum_i|z_i|^2$ without the constraint:
for each nonzero shape, the optimal scale has quadratic moment $n$.
Let $B$ have off-diagonal entries $-(z_i-z_j)^{-2}$ and diagonal
entries $\sum_{j\ne i}(z_i-z_j)^{-2}$.
The quadratic variation is
$-\operatorname{Re}(h^TBh)-\lambda\|h\|^2$.
Varying the phase of $h$ gives $|h^TBh|\le\lambda\|h\|^2$;
Takagi factorization of the complex symmetric matrix implies
$\|B\|\le\lambda$.
Since $B\mathbf1=0$ and $B^T=B$, projecting both coordinate vectors
onto $\mathbf1^\perp$ gives
\[
 |B_{ij}|\le\lambda\left(1-\frac1n\right)
             =\frac{(n-1)^2}{2n}.
\]
This proves the distance bound.
\end{proof}

For a positive integer $q$, put $R=\lceil q\sqrt n\rceil$ and consider
all $n$-element subsets of
\[
 \mathcal G_{n,q}=\{a+ib:\ a,b\in\mathbb Z,\ |a|,|b|\le R\}.
\]
For each such subset $v=(v_1,\ldots,v_n)$, define
\[
 D(v)=n\sum_i|v_i|^2-\left|\sum_iv_i\right|^2,\qquad
 z(v)_i=\frac{n(v_i-\mu_v)}{\sqrt{D(v)}},\qquad
 \mu_v=\frac1n\sum_i v_i.
\]
As in the lattice construction, these points have moment $n$ and their product is
rational. Let
\begin{equation}\label{eq:grid-maximum}
 B_n(q)=\max_{\substack{v\subset\mathcal G_{n,q}\\|v|=n}}
 \left(\frac{n^2}{D(v)}\right)^N
 \prod_{i<j}|v_i-v_j|^2.
\end{equation}
The maximum in this finite set can be found by integer comparisons.

\begin{theorem}\label{thm:grid-global}
For every integer $q\ge2(n-1)$, with
$K_n=n(n-1)(2n-1)/4$, one has
\begin{equation}\label{eq:grid-global}
 B_n(q)\le M_n\le B_n(q)\exp(K_n/q^2).
\end{equation}
If also $q^2>K_n$, the entirely rational enclosure
\begin{equation}\label{eq:grid-rational}
 B_n(q)\le M_n\le\frac{B_n(q)}{1-K_n/q^2}
\end{equation}
holds. In particular, for $0<\varepsilon<1$, choosing
$q\ge\max\{2(n-1),\sqrt{K_n/\varepsilon}\}$ constructs a configuration
with product at least $e^{-\varepsilon}M_n$.
\end{theorem}
\begin{proof}
Fix a centered global maximizer $z^*$. Round the real and imaginary
parts of each $qz_i^*$ to nearest integers, writing $v_i/q=z_i^*+e_i$.
Because $|z_i^*|\le\sqrt n$, this lies in the stated grid.
Let $c=e-n^{-1}(\sum_i e_i)\mathbf1$ and $\rho^2=\|c\|^2$.
Then
\[
 \rho^2\le\frac{n}{2q^2},\qquad
 v_{ij}:=\frac{c_i-c_j}{z_i^*-z_j^*},\qquad
 |v_{ij}|\le b:=\frac{n-1}{\sqrt n\,q}\le\frac12.
\]
Thus the rounded points are distinct. Since $c$ is centered,
$\sum_{i<j}|c_i-c_j|^2=n\rho^2$, so
\[
 \sum_{i<j}|v_{ij}|^2\le\frac{(n-1)^2}{2}\rho^2.
\]
The convergent logarithm series gives
\[
 2\operatorname{Re}\log(1+v)\ge
 2\operatorname{Re}v-\frac{|v|^2}{1-b}\quad (|v|\le b<1).
\]
Put $a=\operatorname{Re}\sum_i\overline z_i^*c_i$.
Stationarity yields $2\operatorname{Re}\sum_{i<j}v_{ij}=(n-1)a$.
The moment before normalization is $n+2a+\rho^2$.
Using $\log(1+x)\le x$ for $x>-1$ in the normalization term, the
linear contributions cancel exactly, and hence
\begin{align*}
 \log\Delta_n(z(v))-\log M_n
 &\ge-\left(\frac{n-1}{2}+
       \frac{(n-1)^2}{2(1-b)}\right)\rho^2\\
 &\ge-\left(\frac{n-1}{2}+(n-1)^2\right)\frac{n}{2q^2}
   =-\frac{K_n}{q^2}.
\end{align*}
The finite grid maximum is at least this rounded value; its own
configuration is feasible, giving both sides of
\eqref{eq:grid-global}. For $0\le x<1$, $e^x\le(1-x)^{-1}$,
which proves \eqref{eq:grid-rational}.
\end{proof}

The guaranteed error concerns a \emph{complete} grid maximum.
Using only a searched subset provides a lower bound and cannot be
substituted for $B_n(q)$ in the upper bound. With the above choice of
$q$, the integer coordinates have height $O(n^2/\sqrt\varepsilon)$,
but the number of subsets is exponential in $n$. High-accuracy
six-point grid enumeration has not been carried out.
The standard-library checker completely enumerates $1176$ subsets
for $n=2,q=2$, giving $B_2(2)=4$, and $1\,873\,200$ subsets for
$n=3,q=4$, giving
\begin{equation}\label{eq:grid-three}
 B_3(4)=\frac{49280427}{1825346}<27.
\end{equation}
An attaining integer triple is $(-7-7i,-7+7i,5)$, with $D=582$.
The corresponding grid enclosure has logarithmic gap at most $15/32$.

\section{An exact polynomial test for stationary configurations}
\label{sec:trace-stationary}
Encoding an unordered configuration by its root polynomial removes
label permutations. A separate positivity condition is needed to
ensure that the force values are paired with the correct conjugate
roots, rather than merely permuted conjugate roots.

\begin{theorem}\label{thm:trace-stationary}
Let $P(t)=t^n+c_2t^{n-2}+\cdots+c_n$ be squarefree, and put
$\mathcal A=\C[t]/(P)$. Let $R$ be the unique polynomial of degree
less than $n$ satisfying
\begin{equation}\label{eq:force-modulo}
 P''-2RP'\equiv0\pmod P.
\end{equation}
For $0\le j,k<n$, define
\begin{equation}\label{eq:trace-matrix}
 H_{jk}=\operatorname{Tr}_{\mathcal A}(t^jR^k).
\end{equation}
The roots $w_i$ of $P$ satisfy
$\sum_{\ell\ne i}(w_i-w_\ell)^{-1}=\overline w_i$ if and only if
\begin{equation}\label{eq:trace-conditions}
 \overline P(R)\equiv0\pmod P,\qquad H=H^*,\qquad H\succ0.
\end{equation}
Here $\overline P$ means coefficientwise conjugation.
In this case $\sum_i|w_i|^2=N$ and the normalized product is
\begin{equation}\label{eq:disc-normalization}
 \Delta_n\!\left(\sqrt{\frac2{n-1}}\,w\right)
   =\left(\frac2{n-1}\right)^N|\operatorname{disc}P|.
\end{equation}
\end{theorem}
\begin{proof}
Squarefreeness makes $P'$ invertible modulo $P$, proving existence
and uniqueness of $R$. At a root,
\[
 R(w_i)=\frac{P''(w_i)}{2P'(w_i)}
       =\sum_{\ell\ne i}\frac1{w_i-w_\ell}.
\]
If this equals $\overline w_i$, both congruence and positivity hold:
the trace is the sum over the roots, so $H=VV^*$ with
$V_{ji}=w_i^j$, an invertible Vandermonde matrix.

Conversely, the congruence says each $R(w_i)$ is a conjugate root.
Positive definiteness makes $H$ nonsingular. Since
$H_{jk}=\sum_i w_i^jR(w_i)^k$, this forces the $R(w_i)$ to be distinct.
Thus $R(w_i)=\overline w_{\sigma(i)}$ for a permutation $\sigma$, and
$H=V\Pi_\sigma V^*$. Consequently $\Pi_\sigma$ is Hermitian positive
definite. A Hermitian permutation is an involution; every
nontrivial transposition has eigenvalue $-1$. It follows that
$\sigma$ is the identity, proving the actual conjugate pairing.
Finally $\sum_i w_iR(w_i)=N$ by pairwise summation.
The discriminant is $\prod_{i<j}(w_i-w_j)^2$, whose absolute value,
not squared absolute value, equals $\Delta_n(w)$.
\end{proof}

This is an exact equivalence, not merely a necessary system.
Let $T$ be the companion matrix representing multiplication by $t$.
The entries $H_{jk}=\operatorname{tr}(T^jR(T)^k)$ are polynomial in
coefficient variables. Taking $R$ as an auxiliary unknown makes both
remainders polynomial. Positive definiteness already enforces
squarefreeness: for distinct roots $\xi_1,\ldots,\xi_r$ with
multiplicities $m_i$, quotient-algebra traces give
$H_{jk}=\sum_{i=1}^r m_i\xi_i^jR(\xi_i)^k$, of rank at most $r$.
Thus $H\succ0$ forces $r=n$.
Hermitian symmetry and positive leading principal minors give real
polynomial conditions, without a separate discriminant denominator.

\begin{corollary}\label{cor:trace-global}
Allow $P$ and the auxiliary polynomial $R$ to vary subject to
\eqref{eq:force-modulo} and \eqref{eq:trace-conditions}. Then
\begin{equation}\label{eq:trace-global}
 M_n=\left(\frac2{n-1}\right)^N\max\det H.
\end{equation}
\end{corollary}
\begin{proof}
Every global maximizer gives a stationary configuration after scaling
by $\sqrt{(n-1)/2}$, and every permitted pair gives a genuine
stationary configuration of moment $N$. For these configurations,
$H=VV^*$, hence $\det H=|\det V|^2=\Delta_n(w)$.
\end{proof}
This is an exact real-polynomial reduction of the global problem;
maximizing its objective has not been carried out for six points.

\begin{proposition}[A genuine spurious six-root pairing]
Let $p$ be the unique root in $(11/25,9/20)$ of
$8p^3-60p^2+126p-45=0$, and put
\[
 s=2p^2-10p+\frac{15}{2},\qquad C=\frac{15}{2}-2p,
 \qquad P(t)=(t^4-st^2+p^2)(t^2-C).
\]
The polynomial
\[
 R(t)=\left(p^2-6p+\frac{41}{4}\right)t
   +\left(-\frac{32p^2}{45}+4p-\frac{26}{5}\right)t^3
   +\left(\frac{29}{75}-\frac{4p^2}{225}\right)t^5
\]
satisfies both force and conjugate-root congruences, and its trace
matrix is Hermitian, but is not positive definite.
\end{proposition}
\begin{proof}
The cubic changes sign on the given rational interval and its
 derivative is positive there. Exact bounds give
$17/5<s<347/100$ and $C>33/5$.
If $A<B$ are the roots of $a^2-sa+p^2$, then $0<A<B<C$.
Thus the six roots are $\pm\sqrt A,\pm\sqrt B,\pm\sqrt C$.
Reduction by the cubic verifies
\begin{align*}
 P''-2RP'&\equiv0\pmod P,&P(R)&\equiv0\pmod P,\\
 tR&\equiv p\pmod{t^4-st^2+p^2},&
 tR&\equiv C\pmod{t^2-C}.
\end{align*}
Since $AB=p^2$ and $p>0$, the force map swaps the positive roots
$\sqrt A,\sqrt B$ and their negative counterparts, while fixing
$\pm\sqrt C$. Its permutation is an involution, so $H$ is Hermitian.
Its principal minor with indices $0,2,4$ equals
\[
 -8(s^2-4p^2)(p^2-Cs+C^2)^2
       =-48600+19440p+2592p^2<0.
\]
Consequently $H$ is indefinite, with inertia $(4,2,0)$.
All displayed identities are independently verified in
$\Q[p]/(8p^3-60p^2+126p-45)$ by the ancillary checker.
\end{proof}

\begin{proposition}[Classical collinear benchmark]\label{prop:collinear}
Among all collinear configurations satisfying the original moment
constraint, the exact maximum is
\begin{equation}\label{eq:collinear}
 R_n^{\rm line}=\frac{\prod_{k=1}^n k^k}{(n-1)^N}.
\end{equation}
In particular, $R_6^{\rm line}=(2/729)C_6$, so every collinear
six-point configuration has product strictly below the central pentagon.
\end{proposition}
\begin{proof}
Center and rotate the line to the real axis. In each ordered chamber,
$\log\Delta_n(x)-\sum_i x_i^2$ is strictly concave, tends to $-\infty$
at collisions and infinity, and has a unique maximizer. Its root
polynomial $P_n$ satisfies
\[
 P_n''-2tP_n'+2nP_n=0,\qquad
 P_n'=nP_{n-1},\qquad
 P_n=tP_{n-1}-\frac{n-1}{2}P_{n-2}.
\]
These are the monic physicists' Hermite polynomials, with $P_0=1$,
$P_1=t$. Pairwise summation gives moment $N$ at this scale.
The recurrence for the absolute resultant is
\[
 |\operatorname{Res}(P_n,P_{n-1})|
 =\left(\frac{n-1}{2}\right)^{n-1}
     |\operatorname{Res}(P_{n-1},P_{n-2})|.
\]
Since $P_n'=nP_{n-1}$, it follows that
$|\operatorname{disc}P_n|=2^{-N}\prod_{k=1}^n k^k$.
Rescaling the moment from $N$ to $n$ gives \eqref{eq:collinear}.
\end{proof}

\section{A finite terminating exact-value algorithm}
\label{sec:exact-algorithm}
A second reduction gives low-degree critical equations directly.
For any configuration choose a farthest pair, apply a similarity
sending it to $w_1=-1,w_2=1$, and write
$w_i=x_i+iy_i$ for $3\le i\le n$.
All remaining points satisfy $|w_i-1|,|w_i+1|\le2$, hence
$|x_i|\le1$, $|y_i|\le\sqrt3$. With
\[
 \mu=\frac1n\sum_iw_i,\qquad Q=\sum_i|w_i-\mu|^2\ge2,
\]
the normalized product is the rational function
\begin{equation}\label{eq:farthest-chart}
 F(w)=\frac{n^N\prod_{i<j}|w_i-w_j|^2}{Q^N}.
\end{equation}
Every distinct real configuration in this chart, even outside the
farthest-pair lens, normalizes to a feasible configuration. Therefore
its unrestricted real maximum is exactly $M_n$, in $2n-4$ variables.
No reflection or polygon symmetry is assumed.

For complexification let $Z_i,W_i$ be independent variables,
$Z_1=W_1=-1$, $Z_2=W_2=1$, and put
\[
 \mu_Z=\frac1n\sum_iZ_i,\quad \mu_W=\frac1n\sum_iW_i,\quad
 Q=\sum_iZ_iW_i-n\mu_Z\mu_W,\quad
 V(Z)=\prod_{i<j}(Z_i-Z_j).
\]
Define, for $i\ge3$,
\begin{equation}\label{eq:critical-low-degree}
 G_i=Q\sum_{j\ne i}\prod_{k\ne i,j}(Z_i-Z_k)
       -N(W_i-\mu_W)\prod_{j\ne i}(Z_i-Z_j),
\end{equation}
and define $J_i$ by swapping $Z$ and $W$.
There are $2n-4$ equations, each of total degree at most $n$.
On the set $QV(Z)V(W)\ne0$, they are exactly the critical equations
of $F=n^NV(Z)V(W)/Q^N$.

\begin{theorem}\label{thm:exact-algorithm}
There is a nonzero polynomial $p_n\in\Q[T]$, computable from
\eqref{eq:critical-low-degree} by saturated elimination, such that
\begin{equation}\label{eq:degree-bound}
 p_n(M_n)=0,\qquad \deg p_n\le n^{2n-4}.
\end{equation}
Combining real-root isolation for $p_n$ with the rational enclosures
\eqref{eq:grid-rational} gives a finite terminating algorithm to
identify the exact value $M_n$ for every fixed $n$.
This is an algorithmic characterization; $p_n$ has not been computed
here for $n\ge4$.
\end{theorem}
\begin{proof}
The real maximum in the chart is a critical point, since it is also
a maximum over the entire real chart and has no collisions.
Saturate the ideal generated by $G_i,J_i$ at $QV(Z)V(W)$, and impose
$TQ^N-n^NV(Z)V(W)=0$.
On each irreducible component meeting that nonvanishing set, all
partial derivatives of $F$ vanish. In characteristic zero, a rational
function with zero differential on an irreducible variety is constant.
Thus there are finitely many complex critical values, even if some
critical components have positive dimension. Elimination to $\Q[T]$
gives a nonzero polynomial vanishing on them and on $M_n$.

Repeated affine B\'ezout bounds give at most $n^{2n-4}$ irreducible
components: intersect successively with the $2n-4$ hypersurfaces of
degree at most $n$; the sum of the degrees of retained components
increases by at most a factor $n$ at each step. Saturation removes
components. There are no more distinct critical values than components.
Their set is invariant under the rational coefficient field, so its
squarefree product polynomial is rational and has the stated degree.

Now isolate the distinct real roots of $p_n$ exactly. For increasing
$q$ with $q\ge2(n-1)$ and $q^2>K_n$, completely compute $B_n(q)$.
The rational interval in \eqref{eq:grid-rational} contains $M_n$ and
has length at most $M_n(K_n/q^2)/(1-K_n/q^2)$, which tends to zero.
Eventually it contains exactly one distinct real root of $p_n$.
That root is $M_n$, irrespective of whether other roots arise from
nonphysical complex components. All enumerations and algebraic
operations at each finite stage terminate. Once the value is isolated,
real-algebraic sampling in its maximizing fiber supplies algebraic
coordinates attaining it \cite{Basu}.
\end{proof}

For $n=2$ there are no free variables and $p_2(T)=T-4$.
For $n=3$ write $Z_3=u,W_3=v$; then
\[
 Q=2+\frac23uv,\qquad
 F=\frac{108(1-u^2)(1-v^2)}{Q^3}.
\]
The two critical equations reduce to
\[
 6u-u^2v+3v=0,\qquad6v-v^2u+3u=0.
\]
Multiplying respectively by $v$ and $u$ and subtracting gives
$v^2=u^2$. The case $v=u$ gives $u=0,\pm3$ and value $27/2$;
the case $v=-u$ gives $u=0,\pm i\sqrt3$ and values $27/2,27$.
Consequently
\[
 p_3(T)=(2T-27)(T-27),\qquad M_3=27.
\]
The fully enumerated lower bound \eqref{eq:grid-three} exceeds $27/2$,
so the grid enclosure isolates the larger critical value.
It is attained by normalizing the triangle $(-1,1,i\sqrt3)$.
For six points the chart has eight real variables and the degree
bound is $6^8=1\,679\,616$; this bound is not a computed six-point
polynomial, nor a claim that such elimination is inexpensive.

\section{What would complete a finite-point global construction?}
The two distinct requirements are to produce a configuration of
value $D$ and to prove $\Delta_n(z)\le D$ for every admissible $z$.
The bounds and candidates above achieve the first requirement and
give a global comparison interval; they do not close that interval
for all finite $n$.
Theorem~\ref{thm:six-basin} now certifies a concrete neighborhood of
the proposed six-point maximum. A global six-point proof still needs
to control the complement. Theorem~\ref{thm:ten-local-trap} demonstrates
why local maximality is insufficient, even at a highly symmetric
candidate.

Theorem~\ref{thm:exact-algorithm} supplies a specific finite algorithm,
with low-degree equations and a degree bound, for every fixed point
count. To determine $M_6$, the remaining task is to compute its full
critical-value eliminant and isolate the physical maximum, or to
solve the coefficient system in Section~\ref{sec:trace-stationary}
and compare $\det H$ on every permitted branch.
The exact spurious pairing in that section shows why genuine
conjugate matching must be enforced. Candidate-adapted nonnegative
certificates could provide a shorter alternative to full elimination.
Permutation-invariant polynomial reduction is established in vortex
computations \cite{FS}; its exhaustive six-point algebraic list and
certified product comparisons have not been reproduced here.

\section{Reproducibility and interpretation}
The standard-library program \texttt{verify\_constructions.py} computes
the normalized integer-lattice products exactly, evaluates the
rational bounds, and checks the ten-point comparison.
It uses explicit exceptions, which remain active under Python
\texttt{-O}. Run it with
\begin{verbatim}
python verify_constructions.py --max-n 30
python verify_hermite.py
python verify_saved_candidates.py
python verify_polygon_perturbations.py
python verify_global_grid.py
python verify_stationary_polynomial.py
python verify_critical_value.py
\end{verbatim}
The optional numerical search is reproduced separately with
\texttt{python explore\_candidates.py}; it requires NumPy.
The JSON file separates exact fractions from floating-point display
values. Finite arithmetic checks support the formulas; the Gaussian,
interpolation and limiting arguments are proved analytically in this note.
The new perturbation checker uses exact cyclotomic field arithmetic
for the ten-point constants and independent root-of-unity power sums,
and verifies the rational six-point neighborhood constants.
The global-grid checker completely enumerates the two stated grids.
The stationary-polynomial checker verifies force remainders and trace
matrices, including an exact spurious six-root pairing. The critical-value
checker verifies the complete three-point elimination identities and
the Hermite resultant recurrence. An optional SageMath eliminant
generator is included as \texttt{critical\_value\_eliminant.sage.py};
it has not been run in this environment and is not a six-point
completeness certificate.
Neither a successful check nor a numerical search is a certificate
that a finite candidate is the unrestricted global maximizer.

\section*{Acknowledgement of assistance}
AI assistance through OpenAI's Codex was used for mathematical
exploration, verification support and manuscript preparation.
The author is responsible for checking the final claims and sources.

\begin{thebibliography}{9}
\small
\setlength{\itemsep}{2pt}
\bibitem{Chalk}
J.~H.~H. Chalk, \emph{A note on a discriminant inequality},
Canadian Mathematical Bulletin \textbf{3} (1960), no.~1, 7--9.
\href{https://doi.org/10.4153/CMB-1960-002-x}{doi:10.4153/CMB-1960-002-x}.
\bibitem{Wang}
J.~Wang, \emph{The Five-Point Case of Mordell's Discriminant Inequality},
arXiv:2610.00358v1 (2026). \url{https://arxiv.org/abs/2610.00358}.
\bibitem{AOC}
Y.~Ameur and J.~Ortega-Cerd\`a,
\emph{Beurling--Landau densities of weighted Fekete sets and correlation
kernel estimates}, Journal of Functional Analysis \textbf{263}
(2012), no.~7, 1825--1861. \url{https://arxiv.org/abs/1110.0284}.
\bibitem{SS}
E.~Sandier and S.~Serfaty, \emph{2D Coulomb Gases and the Renormalized
Energy}, arXiv:1201.3503v2 (2013).
\url{https://arxiv.org/abs/1201.3503}.
\bibitem{Basu}
S.~Basu, \emph{Algorithms in Real Algebraic Geometry: A Survey},
arXiv:1409.1534 (2014). \url{https://arxiv.org/abs/1409.1534}.
\bibitem{CS}
H.~E. Cabral and D.~S. Schmidt,
\emph{Stability of relative equilibria in the problem of $N+1$ vortices},
SIAM Journal on Mathematical Analysis \textbf{31} (1999), no.~2, 231--250.
\href{https://doi.org/10.1137/S0036141098302124}{doi:10.1137/S0036141098302124}.
\bibitem{FS}
J.-C. Faug\`ere and J.~Svartz,
\emph{Solving polynomial systems globally invariant under an action of
the symmetric group and application to the equilibria of $N$ vortices
in the plane}, Proceedings of ISSAC 2012, 170--178.
\href{https://doi.org/10.1145/2442829.2442856}{doi:10.1145/2442829.2442856}.
\end{thebibliography}
\end{document}
